% Unit 17 — termwise geometric series
\input{technique-note.tex}
\title{Integration Bee --- Unit 17:\\ Termwise geometric series}
\author{UK Uni Integration Bee prep\\ Close the sum, then integrate}
\date{}
\begin{document}
\maketitle
\section*{Status}
Tags: \texttt{yang:2.8}, \texttt{infinite-functions}.
A geometric series on a compact subinterval of $[0,1)$ may be integrated term by term. Yang \S2.8 is the contest packaging. The speed move is to sum first. Unit 08 is the first finite truncation. Unit 32 returns to the same geometric sum on $[0,1]$.
\section{Official anchors}
\begin{itemize}
  \item $\sum_{n=0}^{\infty}x^n=1/(1-x)$ for $|x|<1$, and $\sum_{n=1}^{\infty}n x^{n-1}=1/(1-x)^2$.
  \item Syllabus §1 series.
  \item Yang Topics Guide 2nd Ed §2.8 Infinite Functions.
\end{itemize}
\section{Learning goals}
\begin{enumerate}
  \item Recognise a geometric kernel and replace the sum by its closed form before integrating.
  \item Differentiate the geometric sum when the coefficient is $n$.
  \item Integrate a finite truncation directly. Do not open an infinite series for three terms.
\end{enumerate}

\section{Worked geometric sum --- \texttt{U17-D01}}
On $[0,1/2]$,
\[
\sum_{n=0}^{\infty}x^n=\frac{1}{1-x},
\]
so
\[
\int_0^{1/2}\sum_{n=0}^{\infty}x^n\,dx=\int_0^{1/2}\frac{dx}{1-x}=\ln 2.
\]
That is the stored answer. The endpoint $1/2$ is safely inside the disk. \texttt{U17-D08} is the same kernel on the shorter interval $[0,1/3]$:
\[
\int_0^{1/3}\frac{dx}{1-x}=\ln\frac32.
\]
\texttt{U17-D03} is $\int_0^{1/2}dx/(1-x^2)=(\ln 3)/2$, the even cousin, summed as a difference of geometrics or integrated directly.

\section{Worked derivative of the sum --- \texttt{U17-D07}}
Differentiate $\sum_{n=0}^{\infty}x^n=1/(1-x)$ for $|x|<1$:
\[
\sum_{n=1}^{\infty}n x^{n-1}=\frac{1}{(1-x)^2}.
\]
Integrate on $[0,1/2]$:
\[
\int_0^{1/2}\frac{dx}{(1-x)^2}=\frac{1}{1-x}\Big|_0^{1/2}=1.
\]
That is the stored answer. Termwise, $\int_0^{1/2} n x^{n-1}\,dx=(1/2)^n$, and $\sum_{n=1}^{\infty}(1/2)^n=1$. Either path is the check. Use the closed form in a contest.

\section{Faded --- stop early}
\texttt{U17-D09} is the truncation $1+x+x^2$ on $[0,1]$:
\[
\int_0^1(1+x+x^2)\,dx=1+\frac12+\frac13=\frac{11}{6}.
\]
Three terms are not an infinite sum. Integrate them. \texttt{U17-D04} is $\int_0^1 e^x\,dx=e-1$, the Taylor series of $e^x$ summed first. \texttt{U17-D06} is the same integral with the constant term removed, $e-2$.

\section{One substitution}
\texttt{U17-D02} is $\int_0^{\infty}e^{-x}/(1+e^{-x})\,dx=\ln 2$. Let $u=e^{-x}$. \texttt{U17-D05} splits into a power and a square root. A series is the long way around either one.

\section{Drill pointers}
See \texttt{drills.tex} / \texttt{drills.md}. Feed: \texttt{verify/drills\_unit17.json}.

\section*{On a miss}
Teach this file. Write the closed form of the series, then integrate that. If only finitely many powers are present, integrate those powers. Do not invent an answer that is not in \texttt{drills.md}.
\end{document}
